\documentclass[11pt]{article}

% write-latex compliant: no fontspec, no XeLaTeX-only packages.
\usepackage[utf8]{inputenc}
\usepackage[T1]{fontenc}
\usepackage{lmodern}
\usepackage[margin=1in]{geometry}
\usepackage{amsmath, amssymb, amsthm, mathtools}
\usepackage{longtable}
\usepackage{booktabs}
\usepackage{enumitem}
\usepackage{hyperref}
\hypersetup{colorlinks=true, linkcolor=blue, urlcolor=blue, citecolor=blue}

\title{Tour --- {{paper_title}}}
\author{ {{first_author}} et al., {{year}} \\ \small Study repo: \href{ {{repo_url}} }{ {{github_username}}/{{slug}} } }
\date{}

\begin{document}
\maketitle

\noindent
A self-contained walking tour of this study. Read top-to-bottom for a complete first pass; jump by section if you already know the foundations. Estimated total time: \textbf{ {{tour_total_minutes}} minutes}.

\iffalse
NOTE: this file is generated by the study-paper skill, Stage 7.
Mustache-style placeholders {{...}} are replaced per-paper before compilation.
Comments belong inside this iffalse block to keep them out of the rendered PDF.
\fi

\section*{1. Reader's contract}

By the end of this tour you will be able to:
\begin{itemize}[leftmargin=*]
  \item {{learning_outcome_1}}
  \item {{learning_outcome_2}}
  \item {{learning_outcome_3}}
\end{itemize}

\noindent
This tour assumes comfort with {{prereq_topic_1}} (refresh: \href{ {{prereq_link_1}} }{link}) and {{prereq_topic_2}} (refresh: \href{ {{prereq_link_2}} }{link}).

\noindent
Out of scope: {{out_of_scope_1}}; {{out_of_scope_2}}.

\section*{2. Foundations walk}

Each row is one foundation node. Links resolve to the canonical concept page in the shared \href{https://github.com/pleyva2004/math-foundations}{math-foundations} repo.

\begin{longtable}{p{0.22\textwidth} p{0.45\textwidth} p{0.25\textwidth}}
\toprule
\textbf{Concept} & \textbf{Why it matters here} & \textbf{Link} \\
\midrule
\endhead
\bottomrule
\endfoot
{{foundations_walk_table_rows}} \\
\end{longtable}

\section*{3. Paper concepts walk}

The paper graph. Read in topological order; each row links to a Markdown concept page (prose plus formal definition) and a Python witness (CPU-runnable, under 30 seconds).

\begin{longtable}{p{0.22\textwidth} p{0.45\textwidth} p{0.25\textwidth}}
\toprule
\textbf{Concept} & \textbf{Role in the paper} & \textbf{Files} \\
\midrule
\endhead
\bottomrule
\endfoot
{{paper_concepts_walk_table_rows}} \\
\end{longtable}

\noindent\textbf{How the pieces fit.} {{paper_narrative_paragraph}}

\section*{4. Improvements walk}

Each proposal in \texttt{05-improvements.tex} gets exactly one validation file:
\begin{itemize}[leftmargin=*]
  \item \textbf{PROOF mode} --- a standalone LaTeX proof under \texttt{proofs/} with theorem statement, complete proof, and discussion of where it would break.
  \item \textbf{MEASUREMENT mode} --- an extended Python prototype under \texttt{improvements/} exposing \texttt{measure() -> dict} that prints a quantitative comparison table; deterministic under a fixed seed.
\end{itemize}

{{improvements_walk_latex}}

\section*{5. What to do next}

You have finished the tour. Pick one:
\begin{enumerate}[leftmargin=*]
  \item \textbf{Reproduce the sandbox claim.} \texttt{cd sandbox \&\& pip install -r requirements.txt \&\& python experiment.py}.
  \item \textbf{Run a measurement.} Pick any improvement marked MEASUREMENT, run \texttt{python improvements/<slug>.py}, change one knob, re-run.
  \item \textbf{Verify a proof.} Pick any improvement marked PROOF, audit the proof in \texttt{proofs/<slug>.tex}, then try to break it via the Discussion section.
  \item \textbf{Form an opinion.} Fill in \texttt{03-opinions.md}.
  \item \textbf{Extend the foundations graph.} Draft the missing concept at \href{https://github.com/pleyva2004/math-foundations}{pleyva2004/math-foundations}.
\end{enumerate}

\end{document}
